%% HYPOTHESIS / EXPLORATORY DOCUMENT\n%% NOTE: Auto-generated exploratory hypotheses. NOT A PROOF.\n\documentclass{article}\n\usepackage{amsmath,amssymb}\n\begin{document}\n\title{Generated hypothesis from explorerfiles.pdf (Hypothesis)}\maketitle\n\section{Context}\nPseudo-entropy: 0.433843\n\begin{verbatim}\n無と時間、空間、そしてものごとの相対性
Masaaki Yamaguchi

すべての数学への大海原へ入られるのは、ナッシュ先生と広中平祐教授の位相空間の全射と単射と全単射
を、次元の下降と移動、次元の射像を、f (x) = x2 + 1, f | : x → y 、全射と全単射、全射と単射は、違い、全
射は、次元の下降は、全体の次元を１次元降ろす単射であり、３次元から１次元に下降する全射、３次元から
２次元へは、全射でも単射でもない。１次元から２次元、３次元へは単射、３次元から３次元の１次元平面と
２次元から２次元へ、１次元から１次元へは、全単射である。２次元から１次元へも全射であり、３次元から
２次元へは、単射でも全射でもない。本当にわかりやすい。私数学者にはなれないのか？。次元の上への射像
は、この２種の移動が上の式であり、全単射は、全射と単射の同値であり、これを、プログラミング言語の

Ruby の式で、簡易に教えてくださったのを、図も入れてくれたのは、ナッシュ先生と広中平祐先生たちであ
り、私が、田村一郎先生の本でつまづいていたのを、察知してくださり、次元の移動の図で、数式と表現だけ
では、わからないのを、実際の講義の威力と同じ表現で伝授されている。この本は、自分に自信がないので、
隠している。専門家になる人は、見ないのかなと思い、がっくりしていないこともないではないが。田村一郎
先生の本で、これが、わかった後に、多様体、そして、ゼータ関数へといけた、基本群が、この先生の本での
大収穫であった。

このプログラミングの線形代数の本は、先生たちの位相空間の次元の計算の思考を、紙へ

の投射としての、そのままの伝授としかおもえない書き方だ。航空機は、３次元平面から１次元平面への全射
での着陸であり、１次元飛ばすので危ないらしい。宇宙人の UFO は、３次元平面から２次元平面への軟着陸
でなり、全射でも単射でも全単射でもなく、トポロジーが停止していて、韻を踏んで、特異点で着陸であり、
安全らしい。種数が３次元平面と２次元平面で変わっていないらしい。宇宙人もよく考えたものだ。
潜在意識下のベータ関数は、顕在意識のガンマ関数で表せられていて、大海原のベータ関数は、誤差関数と
アーベル多様体であり、この情報は、インターネットの記者の人たちに教えてもらえた。未来と過去は、ベー
タ関数であるが、顕在意識のガンマ関数で分かるらしく、ボルツマン分布が、ゼータ関数のエントロピー値の
部品になっていて、このボルツマン分布は、人工知能のエターナルな空間でもある、一般相対性理論の球対称
である平坦な空間における、全空間である、情報空間での、軌跡の軌道を表している。この確率と組み合わせ
論を、富山大学の数学の講義の、笹田耕一先生に積集合と積空間について、質問を講義のすぐ後にしたら、す
ぐ、世間の常識では、数学の条件としての場合分けであり、エターナルな空間では、すべての組み合わせで全
射と単射、全単射があると教えてもらえて、トポロジーの多様体と基本群が、圭一さんのノートと合わせて、
その後の年の２００５年のＳＲＳの参考書と２００７年の勝間さんのフォトリーディングの書き出す方法を学
んで、マインドマップは、思考の脳内地図であり、投射とは違い、数学の集合論と位相空間は、写像と射像で
あり、このマインドマップとは、別の関係であるのは、分かるが、脳の思考の紙への投射といえば、そう取れ
ないとも言えないが、数学のカテゴリー論の射像がわからなっかのが、一気に、トポロジーの多様体と基本
群、ホモロジーとコホモロジー群へと、トポロジーのほとんどの全分野がわかったことであった。行列は、山
口悟先生と小寺平治先生で、これはすぐにわかっていたことである。
一般座標変換で、各グラフの分布と分散が、座標変換の虚数回転での複素行列での、ベータ関数へと、各点

1

と面、体積が移動している。原子の軌道の軌跡が、極性での電気陰性度の偏りから、原子同士の結合と安定、
存在が出来ていて、このミクロレベルがマクロレベルへと、物体に影を表しているらしい。この影が、ベータ
関数の複素行列の回転体として、光を表している物事のことわりを表しているらしい。このことわりが、事物
の現象の理由や、成り立ちが示されている。それで、この影に見える形にマンデルブロの光の放たれるサース
トン・ペレルマン多様体が、プリズムの倍増で、２００３年の数理科学の表紙と同じ、倍増する光の鳥が、姿
を表す壮絶なきれいさの濃淡を見ると、すべての世界の経脈がわかるらしい。広島女学院の枝垂れ桜の通り道
と同じとおもう。通るのが楽しい道である。パターンとして、Jones 多項式の曲線が、多数の結び目の Jones
多項式の絡\end{verbatim}\n\section{Extracted equations}\n\begin{equation}\label{eq:ex0}\nX=2 のゼータ関数がつくるマンデルブロのサーストン・ペレルマン多様体のプリズムの単体の鳥のかけ\n\end{equation}\n\begin{equation}\label{eq:ex1}\nE = mc2 − mv 2 = ||ds2 || = Γ(γ) dxm\n\end{equation}\n\begin{equation}\label{eq:ex2}\nZ
= e−f dV = dvol = e−f + ef ≥ ef − e−f ≥ e−f − ef\n\end{equation}\n\begin{equation}\label{eq:ex3}\nmanifold

= (x, y, z)/Γ\n\end{equation}\n\begin{equation}\label{eq:ex4}\ndxm =

ex log x div(rotE)dxm\n\end{equation}\n\begin{equation}\label{eq:ex5}\ndxm

= 2(cos(ix log x) − i sin(ix log x))\n\end{equation}\n\begin{equation}\label{eq:ex6}\ndx = 2(i sin(ix log x) − cos(ix log x))\n\end{equation}\n\begin{equation}\label{eq:ex7}\ngij = f , F |gij | : x → y, x → y\n\end{equation}\n\begin{equation}\label{eq:ex8}\nx = p log x, f (y) = p log x\n\end{equation}\n\begin{equation}\label{eq:ex9}\ndxm = F f · F (f )\n\end{equation}\n\section{Derived hypothesis equations}\n\begin{equation}\nX=2 のゼータ関数がつくるマンデルブロのサーストン・ペレルマン多様体のプリズムの単体の鳥のかけ \quad (extracted)\\ G(s) := \int_M K(x,s) \Gamma(\alpha(x,s)) \; d\mu(x)\n\end{equation}\n\begin{equation}\nH(x) := (1/N)\sum_{i=1}^N \phi_i(x) + \prod_{i=1}^N \psi_i(x)  \quad (Higgs-like ansatz)\n\end{equation}\n\begin{equation}\nD_s G(s) + \lambda \int_M H(x) K(x,s) d\mu(x) + R_s = 0  (complexity=3)\n\end{equation}\n\begin{equation}\nE = mc2 − mv 2 = ||ds2 || = Γ(γ) dxm \quad (extracted)\\ G(s) := \int_M K(x,s) \Gamma(\alpha(x,s)) \; d\mu(x)\n\end{equation}\n\begin{equation}\nH(x) := (1/N)\sum_{i=1}^N \phi_i(x) + \prod_{i=1}^N \psi_i(x)  \quad (Higgs-like ansatz)\n\end{equation}\n\begin{equation}\nD_s G(s) + \lambda \int_M H(x) K(x,s) d\mu(x) + R_s = 0  (complexity=3)\n\end{equation}\n\begin{equation}\nZ
= e−f dV = dvol = e−f + ef ≥ ef − e−f ≥ e−f − ef \quad (extracted)\\ G(s) := \int_M K(x,s) \Gamma(\alpha(x,s)) \; d\mu(x)\n\end{equation}\n\begin{equation}\nH(x) := (1/N)\sum_{i=1}^N \phi_i(x) + \prod_{i=1}^N \psi_i(x)  \quad (Higgs-like ansatz)\n\end{equation}\n\begin{equation}\nD_s G(s) + \lambda \int_M H(x) K(x,s) d\mu(x) + R_s = 0  (complexity=3)\n\end{equation}\n\begin{equation}\nmanifold

= (x, y, z)/Γ \quad (extracted)\\ G(s) := \int_M K(x,s) \Gamma(\alpha(x,s)) \; d\mu(x)\n\end{equation}\n\begin{equation}\nH(x) := (1/N)\sum_{i=1}^N \phi_i(x) + \prod_{i=1}^N \psi_i(x)  \quad (Higgs-like ansatz)\n\end{equation}\n\begin{equation}\nD_s G(s) + \lambda \int_M H(x) K(x,s) d\mu(x) + R_s = 0  (complexity=3)\n\end{equation}\n\begin{equation}\ndxm =

ex log x div(rotE)dxm \quad (extracted)\\ G(s) := \int_M K(x,s) \Gamma(\alpha(x,s)) \; d\mu(x)\n\end{equation}\n\begin{equation}\nH(x) := (1/N)\sum_{i=1}^N \phi_i(x) + \prod_{i=1}^N \psi_i(x)  \quad (Higgs-like ansatz)\n\end{equation}\n\begin{equation}\nD_s G(s) + \lambda \int_M H(x) K(x,s) d\mu(x) + R_s = 0  (complexity=3)\n\end{equation}\n\begin{equation}\ndxm

= 2(cos(ix log x) − i sin(ix log x)) \quad (extracted)\\ G(s) := \int_M K(x,s) \Gamma(\alpha(x,s)) \; d\mu(x)\n\end{equation}\n\begin{equation}\nH(x) := (1/N)\sum_{i=1}^N \phi_i(x) + \prod_{i=1}^N \psi_i(x)  \quad (Higgs-like ansatz)\n\end{equation}\n\begin{equation}\nD_s G(s) + \lambda \int_M H(x) K(x,s) d\mu(x) + R_s = 0  (complexity=3)\n\end{equation}\n\begin{equation}\ndx = 2(i sin(ix log x) − cos(ix log x)) \quad (extracted)\\ G(s) := \int_M K(x,s) \Gamma(\alpha(x,s)) \; d\mu(x)\n\end{equation}\n\begin{equation}\nH(x) := (1/N)\sum_{i=1}^N \phi_i(x) + \prod_{i=1}^N \psi_i(x)  \quad (Higgs-like ansatz)\n\end{equation}\n\begin{equation}\nD_s G(s) + \lambda \int_M H(x) K(x,s) d\mu(x) + R_s = 0  (complexity=3)\n\end{equation}\n\begin{equation}\ngij = f , F |gij | : x → y, x → y \quad (extracted)\\ G(s) := \int_M K(x,s) \Gamma(\alpha(x,s)) \; d\mu(x)\n\end{equation}\n\begin{equation}\nH(x) := (1/N)\sum_{i=1}^N \phi_i(x) + \prod_{i=1}^N \psi_i(x)  \quad (Higgs-like ansatz)\n\end{equation}\n\begin{equation}\nD_s G(s) + \lambda \int_M H(x) K(x,s) d\mu(x) + R_s = 0  (complexity=3)\n\end{equation}\n\begin{equation}\nx = p log x, f (y) = p log x \quad (extracted)\\ G(s) := \int_M K(x,s) \Gamma(\alpha(x,s)) \; d\mu(x)\n\end{equation}\n\begin{equation}\nH(x) := (1/N)\sum_{i=1}^N \phi_i(x) + \prod_{i=1}^N \psi_i(x)  \quad (Higgs-like ansatz)\n\end{equation}\n\begin{equation}\nD_s G(s) + \lambda \int_M H(x) K(x,s) d\mu(x) + R_s = 0  (complexity=3)\n\end{equation}\n\begin{equation}\ndxm = F f · F (f ) \quad (extracted)\\ G(s) := \int_M K(x,s) \Gamma(\alpha(x,s)) \; d\mu(x)\n\end{equation}\n\begin{equation}\nH(x) := (1/N)\sum_{i=1}^N \phi_i(x) + \prod_{i=1}^N \psi_i(x)  \quad (Higgs-like ansatz)\n\end{equation}\n\begin{equation}\nD_s G(s) + \lambda \int_M H(x) K(x,s) d\mu(x) + R_s = 0  (complexity=3)\n\end{equation}\n\section{Validation plan}\n1) Specify M and measure; discretize.\n2) Propose parametric forms; fit to data/synthetic tests.\n3) Numerical convergence and residual analysis.\n\section{Caution}\nThis is EXPLORATORY/HYPOTHESIS output. Not a proof.\n\end{document}\n